\section*{Chapitre \ref{ch:ll:the-feed-handler} --- Le feed handler}

\begin{solution}{exo:ll:the-feed-handler:1}
Le paquet de la ligne A contient 1\,003--1\,005 et ouvre un trou à 1\,000~; celui de la ligne B apporte 1\,000--1\,002, qui sont publiés aussitôt,
suivis des messages retenus 1\,003--1\,005~: six événements dans l'ordre des numéros de séquence. Le carnet était périmé de l'arrivée du premier paquet à celle du
deuxième, \qty{40}{\micro\second}, et le trou est considéré comme comblé par l'autre ligne.
\end{solution}

\begin{solution}{exo:ll:the-feed-handler:2}
512~KiB (le noyau double la valeur demandée). À raison de 960~octets par datagramme de 64~octets, il en contient 546, comme mesuré.
\end{solution}

\begin{solution}{exo:ll:the-feed-handler:3}
Le prochain instantané arrive en moyenne un demi-intervalle plus tard et, au pire, un intervalle entier~: $250 + 5 = \qty{255}{\milli\second}$ en
moyenne, $500 + 5 = \qty{505}{\milli\second}$ au pire.
\end{solution}

\begin{solution}{exo:ll:the-feed-handler:4}
La plupart des trous sont des messages que l'autre ligne livre quelques microsecondes plus tard, ou qu'un commutateur a réordonnés~: dans la minute du
chapitre, 4\,224 sur 4\,232. Demander tout de suite surchargerait le serveur et le réseau de demandes portant sur des messages déjà en
route. Une attente plus courte envoie davantage de demandes inutiles~; une attente plus longue laisse le carnet périmé plus longtemps chaque fois qu'un message est réellement perdu
sur les deux lignes.
\end{solution}

\begin{solution}{exo:ll:the-feed-handler:5}
$(2 - 1.2) \times 10^6 \times 0.03 = 24\,000$ datagrammes s'accumulent dans la file. Le noyau doit accorder $24\,000 \times 960 \approx 23$~Mo, donc
environ 11.5~Mo doivent être demandés, ce qui exige de relever nettement \texttt{rmem\_max} au-delà des 4~MiB de l'ordinateur portable.
\end{solution}

\begin{solution}{exo:ll:the-feed-handler:6}
L'instantané remplace les messages perdus par l'état auquel ils ont conduit~: après lui, le carnet est égal à celui de l'exécution de référence
à ce numéro de séquence, et chaque événement postérieur est celui de cette exécution. La construction vérifie les deux~: les événements postérieurs au
numéro de séquence de l'instantané sont égaux à ceux de l'exécution de référence, et le carnet final construit à partir de chaque flux est le même.
\end{solution}

\begin{solution}{exo:ll:the-feed-handler:7}
Rassembler les trous compris entre le prochain numéro attendu et le plus grand numéro retenu (les trous présents dans les clés de la table des messages retenus) et les envoyer
dans une seule demande portant sur une liste d'intervalles~; le serveur répond à tous en un aller-retour. Le trou de \qty{2.7}{\milli\second} de la
minute, qui a demandé plusieurs allers-retours, devrait alors se refermer en un seul~: le délai d'attente plus un aller-retour, \qty{0.7}{\milli\second}.
\end{solution}

\begin{solution}{exo:ll:the-feed-handler:8}
La plupart des trous sont comblés par l'autre ligne en quelques microsecondes ou par une demande de retransmission en moins d'une milliseconde~; attendre un instantané
transforme chacun d'eux en carnet périmé pendant au plus un intervalle d'instantané, mille fois plus longtemps~: dans la minute du chapitre, plus
de quatre mille trous auraient chacun attendu un instantané. La cohérence vient de la rétention des messages jusqu'à ce que le trou soit comblé, et non de leur suppression.
\end{solution}

\begin{solution}{pb:ll:the-feed-handler:1}
\begin{enumerate}
\item 500\,000 datagrammes par seconde~; la rafale en demande trois fois plus, soit un taux d'utilisation de 3.
\item $(1.5 - 0.5) \times 10^6 \times 0.02 = 20\,000$.
\item La file se réduit à $500\,000 - 200\,000 = 300\,000$ par seconde~: environ \qty{67}{\milli\second}.
\item Il attend derrière 20\,000 autres datagrammes servis à \qty{2}{\micro\second}~: \qty{40}{\milli\second}.
\item $20\,000 \times 960 = 19.2$~Mo accordés.
\item 9.6~Mo, plus que \texttt{rmem\_max}~: le noyau le plafonne à 4~MiB (8~MiB accordés) sauf si l'administrateur relève la
limite ou si le processus dispose du privilège nécessaire pour la forcer.
\item Le tampon de 8~MiB contient 8\,738 datagrammes~: environ 11\,262 sont perdus.
\item La mémoire que le handler contrôle~: un thread lecteur qui se contente de déplacer les datagrammes de la socket vers un grand anneau (le travail est
fait par un autre thread), ou les files propres de la carte réseau lues directement par \omterm{def:ll:linux-tuning:polling}{contournement du noyau} (kernel bypass, chapitre~13)~; ou un handler plus rapide.
\item Oui, si le serveur le permet~: 11\,262 messages tiennent dans sa fenêtre de 100\,000 messages, en douze demandes d'au plus 1\,000.
\item Un demi-intervalle d'instantané en moyenne, \qty{0.5}{\second}, et un intervalle entier au pire, \qty{1}{\second}, auquel s'ajoute le chargement.
\item À 200\,000 par seconde après la rafale, environ 100\,000 en moyenne et 200\,000 au pire.
\item Cesser de coter l'instrument et annuler les ordres au repos qui dépendent du carnet, jusqu'à ce que l'\omterm{def:ll:the-feed-handler:stale}{indicateur de carnet périmé} soit abaissé.
\item \textbf{Résultat nommé.} Absorber la rafale exige 20\,000 datagrammes de file, 19.2~Mo de tampon de réception accordé (9.6~Mo
demandés, au-delà de la limite de 4~MiB de cette machine)~; lorsque cela échoue et que seul un instantané permet la reprise, le carnet est périmé pendant
\qty{0.5}{\second} en moyenne et \qty{1}{\second} au pire, avec un instantané par seconde.
\item Trois fois plus rapide~: servir 1.5 million de datagrammes par seconde, soit \qty{0.67}{\micro\second} par datagramme.
\item Chaque ligne dispose de sa propre socket et de son propre tampon, si bien que la mémoire double, et le handler doit lire deux fois plus de datagrammes, dont la moitié de
sont des doublons écartés après une comparaison.
\item À 350 millions de messages par seconde dans sa milliseconde la plus chargée, un flux de cette taille ne peut être traité par un seul thread à des dizaines
de nanosecondes par message~: il faut le répartir entre les lignes et les cœurs.
\item La distribution des arrivées par tranche de 10 millisecondes et par milliseconde, tirée des relevés des journées les plus chargées, le temps de service du
handler sous cette charge, et les compteurs de pertes du noyau pour les sockets.
\item Le high-water mark de la file de réception (et toute perte comptée par le noyau).
\item Un intervalle plus court mobilise la bande passante de la plateforme sur un canal que la plupart des abonnés ne lisent jamais~; la firme paie un intervalle plus long en
carnet périmé après chaque trou irrécupérable.
\item Il met en file ce qu'il peut, perd le reste, le signale par son \omterm{def:ll:the-feed-handler:stale}{indicateur de carnet périmé}, et se rétablit par une demande de retransmission ou par le prochain
instantané.
\end{enumerate}
\end{solution}

\begin{solution}{iq:ll:the-feed-handler:1}
Il reçoit les lignes redondantes (depuis la carte réseau, par le noyau ou directement), décode les paquets, arbitre par
numéro de séquence, détecte et comble les trous (autre ligne, retransmission, instantané), normalise chaque message en un enregistrement
maison, le publie sur des anneaux en mémoire partagée avec un \omterm{def:ll:the-feed-handler:stale}{indicateur de carnet périmé}, et tient à jour des compteurs de supervision.
\iqlookfor{toute la chaîne, avec la reprise et un signal de carnet périmé.}
\end{solution}

\begin{solution}{iq:ll:the-feed-handler:2}
Retenir les messages ultérieurs et lever l'\omterm{def:ll:the-feed-handler:stale}{indicateur de carnet périmé}~; attendre brièvement la ligne B~; si le trou n'est pas comblé, demander la retransmission
au serveur pour la plage et appliquer la réponse dans l'ordre~; si le serveur ne peut pas aider, attendre le prochain instantané de l'instrument, le charger,
écarter les messages retenus qu'il reflète déjà, appliquer le reste et baisser l'indicateur.
\iqlookfor{les trois voies de reprise dans l'ordre, et l'indicateur.}
\end{solution}

\begin{solution}{iq:ll:the-feed-handler:3}
À partir de la plus grande rafale attendue~: l'arriéré (taux d'arrivée moins taux de service, multiplié par la durée de la rafale) multiplié par ce que le noyau
facture par datagramme, bien plus que sa charge utile. Le noyau double la valeur fixée et la plafonne à \texttt{rmem\_max}~;
vérifier la valeur accordée avec \texttt{getsockopt} et surveiller les compteurs de pertes.
\iqlookfor{dimensionnement à partir des données, doublement, plafond et coût par datagramme.}
\end{solution}

\begin{solution}{iq:ll:the-feed-handler:4}
Le handler publie son \omterm{def:ll:the-feed-handler:stale}{indicateur de carnet périmé} avec les événements (et les numéros de séquence permettent à un consommateur de voir une remise à zéro sur instantané). Un
consommateur qui tient son propre carnet vérifie l'indicateur avant de l'utiliser, et contrôle des invariants, comme un carnet non croisé.
\iqlookfor{un indicateur explicite propagé avec les données.}
\end{solution}

\begin{solution}{iq:ll:the-feed-handler:5}
Le code aval (carnets, stratégies, journaux) est écrit une seule fois pour chaque plateforme, les enregistrements ont une taille fixe pour les anneaux et les journaux, et
les numéros de séquence les accompagnent. Le coût est une copie par message et des champs que certaines plateformes ne remplissent pas~; les spécificités propres à une
plateforme doivent être portées par l'enregistrement, faute de quoi elles sont perdues.
\iqlookfor{une interface aval unique contre une copie et des détails perdus.}
\end{solution}

\begin{solution}{iq:ll:the-feed-handler:6}
Rejouer des lignes enregistrées en y injectant pertes, réordonnancements, coupures sur une ou sur les deux lignes, et rafales, face à un serveur de
retransmission et à un \omterm{def:ll:protocols-ii-binary-exchange-protocols:channels}{canal d'instantanés} (snapshot) pouvant être mis en échec~; comparer la sortie avec une exécution de référence (événements identiques après
retransmission~; carnet identique et événements ultérieurs après un instantané), dans chaque langage dans lequel le handler est écrit, comme le fait
la construction.
\iqlookfor{rejeu déterministe avec pannes injectées et une référence propre.}
\end{solution}
